\chapter{Measuring MCP Applications}
\label{ch:measuring}

Everything before this chapter was mechanism. Everything after it is
optimisation, and optimisation without measurement is superstition with a build
pipeline attached. So this is the hinge.

The specific difficulty in an agentic system is that the obvious instruments
measure the wrong thing. Requests per second describes your server, and your
server is three percent of what the user waits for. A latency graph averaged
over all calls hides the one that ran during a fan-out and held up a model turn
behind it. This chapter is about which numbers actually move when your system
gets better, and how to collect them without rebuilding anything.

\section{The metrics that matter}

Nine. Not twenty, because a dashboard nobody reads is a dashboard that does not
exist.

\begin{table}[htbp]
\centering
\small
\begin{tabularx}{\textwidth}{@{}lXl@{}}
\toprule
\thead{Metric} & \thead{Why} & \thead{Meridian} \\
\midrule
Time to first token & What the user feels as ``is it working'' & 340\,ms \\
Loop iterations per task & The unit that actually costs money & 3 \\
Round trips per task & Tool calls plus MRTR retries & 3 \\
Loop latency, decomposed & Model, transport, execution, consent & see below \\
Tool-selection accuracy & Catalogue health & eval-only \\
Call success rate & Server health, split by error kind & \\
Context utilisation & Tokens used against the window & 4{,}009 \\
Cache hit rate & Whether you honour \texttt{ttlMs} & 96.2\,\% \\
Cost per task & The number finance will eventually ask about & \$0.0125 \\
\bottomrule
\end{tabularx}
\caption{Nine numbers. If you can only have three, take iterations per task,
cost per task, and cache hit rate.}
\label{tab:metrics}
\end{table}

Two of these are unusual enough to justify.

\textbf{Loop iterations per task} rather than requests per second. Requests per
second is a server metric and it will look wonderful while your users wait eight
seconds, because the expensive part is not the request.

\textbf{Tool-selection accuracy} needs an eval to measure, so most people skip
it. Do not. It is the leading indicator of catalogue rot, and it degrades slowly
enough that nobody notices from the inside.

\section{Decomposing the loop}

A task is a sequence of iterations, and each iteration divides cleanly into four
bands of time. Recording them separately is what turns ``it feels slow'' into a
number with a component's name attached.

\begin{py}
@dataclass
class Iteration:
    """One trip round the loop, with the milliseconds attributed."""
    index: int
    model_ms: float = 0.0
    transport_ms: float = 0.0
    consent_ms: float = 0.0
    tool_calls: int = 0
    parallel: bool = False
    input_tokens: int = 0
    output_tokens: int = 0
    cost_usd: float = 0.0
\end{py}

Four bands, and every millisecond belongs to exactly one. If you cannot say which
band a millisecond went to, that is the bug: not a slow system, an unmeasurable
one.

\bookfig[0.98]{waterfall}{One Meridian task, drawn to scale. Three iterations,
three tool calls, and the shape tells you where to spend your
effort.}{waterfall}

Read that figure like a trace, because that is what it is.

\begin{itemize}[leftmargin=1.6em, itemsep=3pt]
\item Three iterations, and iteration 2 does no tool calls at all. It exists so
the model can read iteration 1's results and say it is finished, and it costs
400\,ms.
\item The two tool bands are 18.0\,ms and 35.9\,ms. Together, 3.9\,\% of the
task.
\item Iteration 1's two calls ran serially here. In parallel they cost 19.5\,ms
instead of 35.9.
\end{itemize}

\keyidea{If your waterfall is mostly blue, you have a round-trip problem, and the
fix is fewer iterations. If your waterfall has a fat green band, you have a
server problem, and the fix is profiling. Most people assume the second and
have the first.}

\subsection{The subtraction that isolates transport}

From Chapter~\ref{ch:transport}, and worth restating as a technique:

\begin{measurebox}{Same call, three transports}
\begin{tabular}{@{}lrr@{}}
\toprule
\thead{Transport} & \thead{Mean} & \thead{Minus control} \\
\midrule
In-process (control) & 0.031\,ms & \\
stdio                & 0.060\,ms & 0.029\,ms \\
Streamable HTTP      & 0.161\,ms & 0.130\,ms \\
\bottomrule
\end{tabular}
\end{measurebox}

The control still serialises and deserialises, so what remains after the
subtraction is transport mechanics and nothing else. That is the only honest way
to answer ``is my server slow or is my transport slow'', and it takes ten minutes
to set up.

\section{Distributed tracing, done properly}

\texttt{traceparent}, \texttt{tracestate}, and \texttt{baggage} are reserved
\texttt{\_meta} keys, carved out as an explicit exception to the reverse-DNS
prefix rule specifically so W3C Trace Context works unchanged.

\begin{wire}
{
  "jsonrpc": "2.0",
  "id": 2,
  "method": "tools/call",
  "params": {
    "name": "assess_account_risk",
    "arguments": { "accountId": "ACC-1042" },
    "_meta": {
      "io.modelcontextprotocol/protocolVersion": "2026-07-28",
      "io.modelcontextprotocol/clientCapabilities": {},
      "traceparent": "00-0af7651916cd43dd8448eb211c80319c-00f067aa0ba902b7-01"
    }
  }
}
\end{wire}

Client side it is one argument:

\begin{py}
body["_meta"] = build_request_meta(
    self.capabilities, self.info,
    protocol_version=self.protocol_version,
    progress_token=progress_token,
    traceparent=traceparent,
)
\end{py}

Server side it lands on the context, ready to become a span parent:

\begin{py}
return RequestContext(
    ...
    traceparent=meta.get(KEY_TRACEPARENT),
    tracestate=meta.get(KEY_TRACESTATE),
    baggage=meta.get(KEY_BAGGAGE),
)
\end{py}

That is the whole integration. Your existing OpenTelemetry backend now shows
host-to-server spans without knowing what MCP is.

\subsection{The span hierarchy worth building}

\begin{plain}
task                                    [ 1,366 ms]
 |- iteration 0                         [   447 ms]
 |   |- model.generate                  [   430 ms]
 |   `- tool.risk.assess_account_risk   [    18 ms]
 |       |- transport.http              [     4 ms]
 |       `- server.execute              [    13 ms]
 |- iteration 1                         [   518 ms]
 |   |- model.generate                  [   482 ms]
 |   |- tool.fraud.screen_account       [    18 ms]
 |   `- tool.marketdata.get_curve       [    17 ms]
 `- iteration 2                         [   400 ms]
     `- model.generate                  [   400 ms]
\end{plain}

Three properties make this worth the effort:

\textbf{Iterations are spans.} Not just tool calls. Otherwise your trace shows
54\,ms of activity inside a 1,366\,ms task and a 1,312\,ms hole nobody can
explain.

\textbf{Transport and execution are separate spans.} That is the subtraction above, done automatically for every call.

\textbf{Sibling spans that overlap in time reveal parallelism.} If you fanned
out, the trace shows it. If you thought you fanned out and did not, the trace
shows that too, which is more useful.

\begin{notebox}{What to put in \texttt{baggage}}
\texttt{baggage} propagates key-value pairs across every hop, which makes it
tempting as a general-purpose channel. Resist.

Good: tenant id, task id, environment. Things you want to slice traces by.

Bad: anything sensitive. Baggage crosses every trust boundary in the system,
including into servers you do not operate, and it lands in whatever tracing
backend each of them uses.
\end{notebox}

\begin{legacybox}{Logging}
Deprecated as of \texttt{2026-07-28} (SEP-2577), eligible for removal no earlier
than 2027-07-28. It still works in the meantime, and it changed shape rather
than disappearing, which is why it is easy to misread as removed.

The \texttt{logging/setLevel} RPC is gone. A client now sets the level it wants
per request, in \texttt{\_meta}:

\begin{wire}
"_meta": { "io.modelcontextprotocol/logLevel": "warning" }
\end{wire}

Absent that key, a server \textbf{must not} send \texttt{notifications/message}
for the request at all. Opting in is the client's move, and it is per request
rather than per connection, for the same reason everything else in this revision
is.

The notifications themselves, the \texttt{logging} server capability, and the
severity levels all survive the deprecation window unchanged.

Migrating: \texttt{stderr} for stdio, where the client already captures it, and
OpenTelemetry for anything structured. The SEP's reasoning is that both are
mature, widely adopted, and better at this than an application protocol is, and
that carrying a second logging system in the core specification taxes every
implementation whether or not it uses one.
\end{legacybox}

\section{Instrumenting a server}

Meridian keeps deliberately crude built-in counters, because something is
enormously better than nothing and it costs four lines:

\begin{py}
finally:
    elapsed = (time.perf_counter() - started) * 1000.0
    self.call_count[method] = self.call_count.get(method, 0) + 1
    self.total_ms[method] = self.total_ms.get(method, 0.0) + elapsed
\end{py}

\begin{py}
def stats(self) -> dict:
    return {
        "calls": dict(self.call_count),
        "meanMs": {
            m: round(self.total_ms[m] / self.call_count[m], 3)
            for m in self.call_count if self.call_count[m]
        },
    }
\end{py}

Means are a lie, of course. What you actually want is percentiles, because the
p99 is what your users complain about and the mean is what hides it. The client
keeps those:

\begin{py}
def percentile(self, p: float) -> float:
    if not self.latency_ms:
        return 0.0
    ordered = sorted(self.latency_ms)
    idx = min(len(ordered) - 1, int(round((p / 100.0) * (len(ordered) - 1))))
    return ordered[idx]
\end{py}

\keyidea{Report p50, p95, and p99. A mean of 40\,ms with a p99 of 4\,s is a
system where one user in a hundred thinks you are broken, and the mean will never
tell you.}

\subsection{Tail amplification, which is why fan-out needs watching}

There is a second-order effect here that agent systems run into harder than
ordinary services, and it follows from the parallel fan-out this book keeps
recommending.

A fan-out finishes when its \emph{slowest} branch finishes. So the task's latency
is the maximum over its branches, and maxima are governed by tails.

Take three independent servers, each with a p99 of two seconds. A task calling
all three in parallel is slow whenever \emph{any} of them is slow. The chance
that none is in its slow 1\,\% is $0.99^3$, about 97\,\%, so roughly 3\,\% of
tasks hit a two-second branch. What was a one-in-a-hundred event at the server
became a one-in-thirty event at the task.

Widen the fan-out and it gets worse quickly:

\begin{table}[htbp]
\centering
\small
\begin{tabular}{@{}rrr@{}}
\toprule
\thead{Parallel calls} & \thead{Tasks hitting a p99 branch} & \thead{Effectively} \\
\midrule
1  & 1.0\,\%  & p99 \\
3  & 3.0\,\%  & p97 \\
10 & 9.6\,\%  & p90 \\
40 & 33.1\,\% & p67 \\
\bottomrule
\end{tabular}
\vspace{6pt}
\footnotesize

Arithmetic, assuming independence: $1 - 0.99^{n}$. Real branches correlate, so
treat this as the optimistic case.
\end{table}

Read the last row against Chapter~\ref{ch:tools}'s advice to add batch tools. A
40-account portfolio review fanned out as 40 parallel calls hits a p99 branch in
a third of all tasks. The same review through \texttt{rank\_portfolio\_risk} is
one call and hits it in one task in a hundred. Batching is a tail-latency
optimisation as much as a token one, and that is the argument nobody makes for
it.

Two practical consequences. Your server's p99 is your task's typical experience
once you fan out past a handful of branches, so it deserves the attention people
reserve for the mean. And a hedge, sending a duplicate request when one branch
passes a threshold, buys more in an agent system than in ordinary RPC, because
the branch you are waiting on is holding up a model turn behind it.

\section{Instrumenting the client}

\begin{py}
@dataclass
class CallStats:
    """The counters Chapter 16 turns into a dashboard."""
    requests: int = 0
    mrtr_rounds: int = 0
    retries: int = 0
    errors: int = 0
    bytes_sent: int = 0
    bytes_received: int = 0
    latency_ms: list[float] = field(default_factory=list)
\end{py}

Three of these are worth more attention than they usually get.

\textbf{\texttt{mrtr\_rounds}} is extra round trips forced by servers asking
questions, at roughly 375\,ms each (Chapter~\ref{ch:mrtr}).
Chapter~\ref{ch:building-hosts} called it the early warning for a missing schema
field; on a dashboard it is also the fastest way to attribute a latency
regression to a specific server deploy.

\textbf{\texttt{bytes\_sent} against \texttt{bytes\_received}} tells you whether
your requests or your responses are the problem. Meridian sends a 320-byte
request and receives 5,856 bytes for a \texttt{tools/list}. Requests are never
the issue; catalogues frequently are.

\textbf{\texttt{errors}, split by kind.} Protocol errors and tool execution
errors have completely different meanings. Protocol errors are bugs. Tool
execution errors are normal operation, and a rate of zero means your validation
is happening somewhere it should not be.

\section{Cache metrics}

\begin{py}
@dataclass
class CacheStats:
    hits: int = 0
    misses: int = 0
    stale: int = 0
    invalidations: int = 0
    stores: int = 0
    uncacheable: int = 0
\end{py}

Six counters, and the split between \texttt{misses} and \texttt{stale} is the
one that pays.

\begin{itemize}[leftmargin=1.6em, itemsep=3pt]
\item High \texttt{misses}: you are seeing new things. Cold cache, or key
cardinality is higher than you thought.
\item High \texttt{stale}: TTLs are too short for your access pattern. This is
tunable and usually free.
\item High \texttt{invalidations}: the underlying data really is changing.
Nothing to fix, but it explains the hit rate.
\item High \texttt{uncacheable}: you are calling uncacheable methods, or making
MRTR retries. Expected for \texttt{tools/call}; suspicious for anything else.
\end{itemize}

A hit rate alone cannot distinguish ``TTL too short'' from ``data genuinely
volatile'', and those have opposite fixes.

\begin{measurebox}{Meridian's cache, 600 reads over 40 URIs}
\begin{tabular}{@{}lr@{}}
\toprule
Hit rate            & 96.2\,\% \\
Requests avoided    & 577 of 600 \\
Wall clock          & 12.4\,ms $\rightarrow$ 1.8\,ms \\
Speedup             & 6.82$\times$ \\
\bottomrule
\end{tabular}
\end{measurebox}

\section{Token and cost accounting}

\begin{py}
def estimate_tokens(text: str) -> int:
    return max(1, round(len(text) / CHARS_PER_TOKEN))

def estimate_cost_usd(input_tokens: int, output_tokens: int) -> float:
    return (input_tokens * INPUT_USD_PER_MTOK / 1_000_000.0
            + output_tokens * OUTPUT_USD_PER_MTOK / 1_000_000.0)
\end{py}

The character-based estimator is wrong in the third decimal place and right
enough to compare a before against an after, which is what it is for. Your
provider reports exact counts; use those in production and the estimator when you
need a number before you have spent the money.

\begin{measurebox}{Meridian's per-iteration cost}
\begin{tabular}{@{}lrrr@{}}
\toprule
\thead{Iteration} & \thead{Input} & \thead{Output} & \thead{Cost} \\
\midrule
0 & 1{,}252 & 12 & \$0.003936 \\
1 & 1{,}318 & 19 & \$0.004239 \\
2 & 1{,}400 & 8  & \$0.004320 \\
\midrule
\textbf{Total} & \textbf{3{,}970} & \textbf{39} & \textbf{\$0.012495} \\
\bottomrule
\end{tabular}
\end{measurebox}

Input climbs, output does not, and iteration 2 is the most expensive despite
doing the least. That is the economics of agent loops in one table.

\subsection{Catalogue cost as a first-class metric}

\begin{py}
def catalogue_tokens(self) -> int:
    return sum(
        estimate_tokens(f"{t.get('name','')}{t.get('description','')}"
                        f"{t.get('inputSchema','')}{t.get('outputSchema','')}")
        for t in self.catalogue())
\end{py}

Put this on a dashboard. It changes whenever anybody adds a tool to any connected
server, it is charged on every turn of every task, and nobody notices it climbing
until it is 12,000 and the model has started guessing.

\section{Evals as regression tests}

The four numbers to track per release:

\begin{table}[htbp]
\centering
\small
\begin{tabularx}{\textwidth}{@{}lXl@{}}
\toprule
\thead{Number} & \thead{Regression means} & \thead{Gate} \\
\midrule
Task success rate       & Something broke behaviourally & no drop $>$ 3 points \\
Iterations per task     & The model needs more turns for the same work & no rise $>$ 10\,\% \\
Tokens per task         & Context or catalogue grew & no rise $>$ 10\,\% \\
Cost per task           & All of the above, monetised & no rise $>$ 10\,\% \\
\bottomrule
\end{tabularx}
\caption{Gate on movement against the last release, not on absolutes. Absolute
thresholds get deleted the first time somebody needs to ship.}
\label{tab:eval-gates}
\end{table}

Catching \emph{capability drift} is the real prize here. The model provider ships
a new version. Your tests pass, because they use a stub. Your evals show
iterations per task moved from 3.0 to 3.4, which is a 13\,\% cost increase nobody
would have seen otherwise.

\section{A dashboard worth having}

\begin{plain}
MERIDIAN                                    last 24h

  TASKS
  completed                     14,204        +2.1%
  p50 wall clock                 1,412 ms     +0.3%
  p95 wall clock                 3,890 ms     +8.2%   <-- look here
  iterations per task (mean)       3.14        +0.1%
  cost per task                  $0.0131      +4.8%

  PROTOCOL
  round trips per task             3.21        +6.9%   <-- and here
  mrtr rounds per task             0.19       +58.0%   <-- and here
  cache hit rate                   94.1%       -2.1pt

  SERVERS                     calls    p95      errors
  risk                       28,401   14 ms      0.02%
  compliance                  9,120   31 ms      0.11%
  fraud                      27,880    8 ms      0.00%
  marketdata                 31,004    3 ms      0.00%
\end{plain}

Read the arrows together and the story writes itself. MRTR rounds up 58\,\%,
round trips up 6.9\,\%, p95 up 8.2\,\%, cost up 4.8\,\%. Somebody deployed a
server that started asking for something it used to infer, and everything
downstream is a consequence.

That diagnosis takes ten seconds with these numbers and a week without them.

\section{Alerting}

Four alerts. Chapter~\ref{ch:operations} adds the operational ones; these are the
protocol ones.

\begin{table}[htbp]
\centering
\small
\begin{tabularx}{\textwidth}{@{}lXX@{}}
\toprule
\thead{Alert} & \thead{Condition} & \thead{Usually means} \\
\midrule
Iteration inflation & Iterations per task up $>$ 20\,\% & Tool or prompt change \\
Catalogue bloat     & Catalogue tokens up $>$ 25\,\% & Somebody added tools \\
Cache collapse      & Hit rate down $>$ 15 points & TTLs changed, or ordering went unstable \\
MRTR surge          & MRTR rounds per task doubled & A server started asking \\
\bottomrule
\end{tabularx}
\caption{All four are ratios against a baseline, because absolute thresholds are
wrong for every deployment except the one that set them.}
\label{tab:alerts}
\end{table}

\section{The harness}

Meridian's measurement rig is about eighty lines, and its shape is worth copying
whatever you end up building it on top of.

\begin{py}
def time_it(fn, *, n: int = 200, warmup: int = 20, label: str = "") -> Timing:
    for _ in range(warmup):
        fn()
    timing = Timing(label)
    for _ in range(n):
        started = time.perf_counter()
        fn()
        timing.add((time.perf_counter() - started) * 1000.0)
    return timing
\end{py}

Warmup matters more than people expect. The first calls pay import costs, JIT
warmup, connection setup, and cold caches, and including them turns a
distribution into a bimodal mess.

Eight scenarios, each answering one question:

\begin{table}[htbp]
\centering
\small
\begin{tabularx}{\textwidth}{@{}lX@{}}
\toprule
\thead{Scenario} & \thead{Question} \\
\midrule
\texttt{transport}     & What does each transport cost, minus execution? \\
\texttt{coldstart}     & Spawn per call, or warm pool? \\
\texttt{catalogue}     & What does catalogue bloat cost per turn? \\
\texttt{cache}         & What is honouring \texttt{ttlMs} worth? \\
\texttt{fanout}        & What does parallelism buy, at what latency? \\
\texttt{mrtr}          & What does one elicitation really cost? \\
\texttt{loop}          & Where does a whole task's time go? \\
\texttt{serialisation} & How much of a request is envelope? \\
\bottomrule
\end{tabularx}
\caption{Each scenario exists to answer a question somebody actually asked.
Benchmarks without a question attached become numbers nobody acts on.}
\label{tab:bench-scenarios}
\end{table}

\begin{sh}
make bench                                   # everything, print a table
python3 -m meridian.bench.run cache fanout   # named scenarios
python3 -m meridian.bench.run --json out.json
\end{sh}

\subsection{Two things a harness must be honest about}

\textbf{State what is simulated.} Meridian's model latency is drawn from a stated
distribution and its docstring says so in the first paragraph. A harness that
quietly models something is generating fiction with error bars.

\textbf{State the machine.} Absolute numbers do not transfer. Ratios do.

\begin{py}
def machine() -> dict:
    return {
        "python": platform.python_version(),
        "platform": platform.platform(),
        "machine": platform.machine(),
        "processor": platform.processor() or platform.machine(),
    }
\end{py}

\section{What to remember}

\begin{itemize}[leftmargin=1.6em, itemsep=3pt]
\item Nine metrics. If you take three: iterations per task, cost per task, cache
hit rate.
\item Decompose every iteration into model, transport, execution, and consent.
Every millisecond belongs to one band.
\item Subtract an in-process control to isolate transport from execution.
\item \texttt{traceparent} in \texttt{\_meta} is one line and gives you real
distributed traces.
\item Make iterations spans, not just tool calls, or your trace has a hole where
96\,\% of the time went.
\item Report percentiles. Means hide the users who are suffering.
\item Split cache misses from staleness. They have opposite fixes.
\item Track catalogue tokens on a dashboard. It grows silently.
\item Gate evals on movement against the last release, not on absolute
thresholds.
\item Warm up before timing, state what is simulated, and state the machine.
\end{itemize}

Next: spending the measurements.
